% Printed 62--71, including the intervening unnumbered leaves.
\clearpage
\noindent\makebox[2em][l]{62}\parbox{\dimexpr\linewidth-4em\relax}{\centering\small NOTICE.\quad DETERMINATION OF THE DIFFERENCE IN LATITUDE ETC.}\hfill\mbox{}\par
\noindent Mr.\ \textsc{Hansen}, made with \textsc{Ertel}'s meridian circle on the same zenith stars
at the same time as the author, by arrangement with him. The result, based on 206 observations,
is of still greater interest because the Seeberg Observatory is connected to the Hanoverian
triangulation network by the triangles measured under the direction of Lieutenant General
von \textsc{Müffling}.

Finally, the work closes with the result of a new calculation of all previous reliable arc
measurements, carried out by Dr.\ \textsc{Schmidt} at the author's request. The principles are similar
to those underlying \textsc{Walbeck}'s work mentioned above, but the calculation is more precise,
takes account of \textls[150]{all} the astronomical points occurring in the individual arc
measurements (rather than only the endpoints), and incorporates the results of the Hanoverian
arc measurement contained in the present work. This determination of the dimensions of the
Earth ellipsoid may be regarded as the most reliable conclusion that can so far be drawn from
all the latitude arc measurements. It is also of interest to see our conception of the degree
of accuracy so far attainable by this means established according to definite principles.
\begin{center}\rule{.25\textwidth}{.4pt}\end{center}

\clearpage
% Unnumbered Bemerkungen opening at the position of printed 63.
\vspace*{14mm}
\begin{center}\textbf{REMARKS.}\end{center}
{\small
In the original printing of the Difference in Latitude, which could also be checked against
\textsc{Gauss}'s surviving manuscript, some slight inaccuracies were found in the numerical
data. These were corrected only in part.

In the foregoing reprint, the original entry on p.\ 9 was changed for
\begin{center}
\par\noindent\begin{tabular}{ccccccl}
the&declination&of&star&26:&$50^\circ28'13''{,}12$&to $50^\circ38'13''{,}12$\\
»&»&»&»&40:&$56\quad42\quad21{,}84$&» $56\quad52\quad21{,}74$,
\end{tabular}
\end{center}
since otherwise the observed values given for these stars on pp.\ 22 and 27 are inconsistent
with the errors remaining after the adjustment on p.\ 35.

In Chapter II, the Observations, small changes of 1 to 2 hundredths of a second were made to
the final values on a few occasions to bring them into agreement with the observed values
and the reductions. This was done on the condition that the mean final values entering into
the subsequent calculation were unaffected. The following entries had to be left uncorrected.
\par\noindent\begin{tabular}{@{}p{.065\linewidth}p{.91\linewidth}@{}}
P.\ 21&For star 25 ($\theta$ Draconis), Altona. Limb west. June 16, the final value obtained
from the observed value and the two reductions is $1''$ greater than stated.\\
P.\ 22&For star 26, Altona. Limb east. June 27, the observed value and the reductions yield
a final value of $\ldots\ 34''{,}98$ instead of $\ldots\ 34''{,}89$. The latter value was used,
however, in taking the mean.\\
P.\ 26&For star 36, Göttingen. Limb west, the mean of the final values is $\ldots\ 21''{,}14$.
The stated value $\ldots\ 21''{,}13$ was used for the calculation.\\
P.\ 27&For star 40, Altona. Limb west. June 22, the final value derived from the observed value
and the reductions is $0''{,}10$ greater than stated there.
\end{tabular}

The following changes were made in Chapter III, Results.
\par\noindent\begin{tabular}{@{}p{.065\linewidth}p{.91\linewidth}@{}}
P.\ 30&The difference in latitude from the observations of star 39, given as $2^\circ0'58''{,}51$
in the original, was changed to $\ldots\ 58''{,}57$, as follows from the values on p.\ 27.
The stated mean value of the difference in latitude on p.\ 30 is unaffected.
The weight of the result is $213{,}42$, as stated. The original gives $213{,}41$ for this.\\
P.\ 34&In Art.\ 6, line 8 from the top, the original value for the correction to the collimation
error in Göttingen, $-0''{,}012$, was changed to $+0''{,}012$, in accordance with the values $3''{,}75$
on p.\ 32, line 2 from the top, and $3''{,}76$ on p.\ 34, line 12 from the top.
Recalculation confirmed that the $+$ sign is correct.
\end{tabular}\par
}

\clearpage
\noindent\makebox[2em][l]{64}\parbox{\dimexpr\linewidth-4em\relax}{\centering\small REMARKS.\quad DETERMINATION OF THE DIFFERENCE IN LATITUDE ETC.}\hfill\mbox{}\par
{\small
\par\noindent\begin{tabular}{@{}p{.065\linewidth}p{.91\linewidth}@{}}
P.\ 45&In the error column of the table
\end{tabular}
\begin{center}
\par\noindent\begin{tabular}{llllll}
the first&value&$-0''{,}31$&in the original&to&$-0''{,}30$,\\
the last&»&$-0{,}23$&»\quad »&»&$-0{,}24$;
\end{tabular}
\end{center}
\begin{quote}
and the sum of the products of the squared errors and the weights, given twice in the original
as $9{,}6184$, was changed to $9{,}6154$.
\end{quote}
The following entry was left unchanged.
\par\noindent\begin{tabular}{@{}p{.065\linewidth}p{.91\linewidth}@{}}
P.\ 31&For star 9, Altona, the weight of the collimation error, $4{,}00$. According to the number
of observations on p.\ 15, it must be $5{,}33$. The value $4{,}00$ was retained because it was used
to derive the final result on p.\ 32.
\end{tabular}

A further remark concerns Art.\ 7, p.\ 35. The sum of the products of the squares of the
differences listed in Art.\ 6 and the corresponding numbers of observations is given there as
$292{,}8249$. Since it did not yield exactly the mean error of an observation derived therefrom,
$1''{,}5308$, these values were recalculated. The recalculation gave $291{,}8422$ in place of the
first value, and hence $\pm1''{,}5280$ as the mean error of an observation and $\pm0''{,}1036$ as the
mean error of the difference in latitude in place of the original $\pm0''{,}1038$.

Finally, in Chapter IV, in the table on p.\ 53 for star 31, the zenith distance stated in the
original, $+4^\circ48'58''{,}28$, was changed to $+4^\circ57'58''{,}28$. The latter value can be
recovered from the latitude given alongside and the declination listed for the star on p.\ 9.

The treatise on the Lapland arc measurement mentioned by \textsc{Gauss} in the Notice of the
Determination of the Difference in Latitude etc., p.\ 61, is probably the one by
O.\ A.\ \textsc{Rosenberger}, Über die, auf Veranstaltung der französischen Academie, während der
Jahre 1736 und 1737 in Schweden vorgenommene Gradmessung. (Astr.\ Nachr.\ Sixth Volume, 1828.
Pp.\ 1--32).
\par\hfill\textsc{Krüger, Börsch.}\par
}
\begin{center}\rule{.25\textwidth}{.4pt}\end{center}

\clearpage
% Unnumbered section title, principal PDF 72 (position 65).
\vspace*{45mm}
\begin{center}
{\Huge EARTH ELLIPSOID}\par\vspace{10mm}
{\large AND}\par\vspace{10mm}
{\Huge GEODESIC.}
\end{center}
\clearpage
% Blank verso, principal PDF 73 (position 66).
\null
\clearpage
% Unnumbered Nachlass opening, principal PDF 74 (position 67).
\begin{center}
{\Large\textls[180]{POSTHUMOUS PAPERS.}}\par\vspace{12mm}
[1.]\par
The [Earth] ellipsoid.
\end{center}
Let $a$ be the semimajor and $b$ the semiminor axis. At a given place let $\varphi$ be the
astronomical latitude, $\varphi'$ the so-called corrected latitude
[and $\psi$ the reduced latitude]. Let $r$ be the distance from the centre, $x$ the
distance from the Earth's axis, and $y$ the distance from the equator.

Then, since the equation of the ellipse gives $\dfrac{xx}{aa}+\dfrac{yy}{bb}=1$, one can set
\[
x=r\cos\varphi'=a\cos\psi,\qquad y=r\sin\varphi'=b\sin\psi
\]
If $s$ denotes the elliptic arc between the place and the equator, one has
\[
-\sin\varphi\cdot\mathrm ds=\mathrm dx,\qquad\cos\varphi\cdot\mathrm ds=\mathrm dy,
\]
and thus
\[
\sin\varphi\cdot\mathrm ds=a\sin\psi\cdot\mathrm d\psi,\qquad\cos\varphi\cdot\mathrm ds=b\cos\psi\cdot\mathrm d\psi.
\]
It follows that
\[
\operatorname{tang}\psi=\frac ba\operatorname{tang}\varphi,\qquad\operatorname{tang}\varphi'=\frac{bb}{aa}\operatorname{tang}\varphi.
\]
Furthermore
\[
\begin{aligned}
\cos\psi&=\frac{a\cos\varphi}{\sqrt{(aa\cos\varphi^2+bb\sin\varphi^2)}}=\sqrt{\frac{\cos\varphi\cos\varphi'}{\cos(\varphi-\varphi')}},\quad&\cos\varphi'&=\frac{aa\cos\varphi}{\sqrt{(a^4\cos\varphi^2+b^4\sin\varphi^2)}}\\
\sin\psi&=\frac{b\sin\varphi}{\sqrt{(aa\cos\varphi^2+bb\sin\varphi^2)}}=\sqrt{\frac{\sin\varphi\sin\varphi'}{\cos(\varphi-\varphi')}},&\sin\varphi'&=\frac{bb\sin\varphi}{\sqrt{(a^4\cos\varphi^2+b^4\sin\varphi^2)}}
\end{aligned}
\]
\[
r=\sqrt{\frac{a^4\cos\varphi^2+b^4\sin\varphi^2}{aa\cos\varphi^2+bb\sin\varphi^2}}=a\sqrt{\frac{\cos\varphi}{\cos\varphi'\cos(\varphi-\varphi')}}=b\sqrt{\frac{\sin\varphi}{\sin\varphi'\cos(\varphi-\varphi')}}
\]
\noindent\hfill $9^*$\par

\clearpage
\noindent 68\hfill POSTHUMOUS PAPERS.\hfill\mbox{}\par
\[
\frac{a\,\mathrm d\psi}{\cos\psi^2}=\frac{b\,\mathrm d\varphi}{\cos\varphi^2},\quad\text{thus}\quad\frac{\mathrm ds}{\mathrm d\varphi}=\frac{bb\cos\psi^3}{a\cos\varphi^3}=\frac{aa\sin\psi^3}{b\sin\varphi^3}=\frac{aabb}{(aa\cos\varphi^2+bb\sin\varphi^2)^{\frac32}}
\]
\[
ee=1-\frac{bb}{aa}.
\]
If one quadrant of the Earth's meridian is $=Q$, then one quadrant of the equator is [cf.\ Volume IV, p.\ 330]
\[
\begin{aligned}
&=Q\cdot(1+\tfrac14ee+\tfrac7{64}e^4+\tfrac{15}{256}e^6+\tfrac{579}{16384}e^8+\tfrac{1515}{65536}e^{10}+\ldots)\\
&=Q:(1-\tfrac14ee-\tfrac3{64}e^4-\tfrac5{256}e^6-\tfrac{175}{16384}e^8-\tfrac{441}{65536}e^{10}-\ldots).
\end{aligned}
\]
\begin{center}\rule{.25\textwidth}{.4pt}\end{center}
For a flattening of $\tfrac1{310}$, $e=\dfrac{\sqrt{619}}{310}$, $\log e=8{,}9044836$
\[
\begin{array}{l@{\quad}l@{\quad}l}
ee\ldots\ 7{,}808\,9672&e^4\ldots\,\ldotp\ 5{,}617\,9344&e^6\ldots\,\ldotp\,\ldotp\ 3{,}426\,9016\\[4pt]
\tfrac14\ldots\,\ldotp\ 9{,}397\,9400&\tfrac7{64}\ldots\ 9{,}038\,9180&\tfrac{15}{256}\ldots\ 8{,}767\,8513
\end{array}
\]
Quadr.\ of the equat.\ [if $Q=10\,000\,000$ metres is assumed:]
\[
\begin{array}{r@{\quad}c@{\quad}r}
10\,016\,103^{\mathrm m}{,}015&&\\
45{,}379&&\\
0{,}157&&\\\cline{1-1}
10\,016\,148^{\mathrm m}{,}551&\ldotp\,\ldotp\,\ldotp\,\ldotp\,\ldotp\,\ldotp\,\ldotp&7{,}000\,7008\\
&\left[\dfrac\pi2\right]\ldots&0{,}196\,1199\\\cline{3-3}
&\log a=&6{,}804\,5809.
\end{array}
\]
\begin{center}\rule{.25\textwidth}{.4pt}\end{center}
\[
\text{Flattening}=\omega,\qquad b=a(1-\omega),\qquad ee=(2-\omega)\omega.
\]
\[
e=\sin E,\qquad\frac{1-\sqrt{(1-ee)}}{1+\sqrt{(1-ee)}}=\operatorname{tang}\tfrac12E^2.
\]
For $\omega=\dfrac1{302{,}68}$ [one has]
\[
E=4^\circ39'31''{,}355,\qquad\operatorname{tang}\tfrac12E\ldots\ 8{,}609\,3534
\]
\[
\begin{array}{l@{\quad}l@{\quad}l}
ee\ldots\ 7{,}819\,3287&e^4\ldots\,\ldotp\ 5{,}638\,6574&e^6\ldots\,\ldotp\ 3{,}457\,9861\\[4pt]
\tfrac14\ldots\ 9{,}397\,9400&\tfrac7{64}\ldots\ 9{,}038\,9180&\tfrac{15}{256}\ldotp\,\ldotp\ 8{,}767\,8513
\end{array}
\]
[For $Q=10000000$ metres the quadrant of the equator is therefore:]

\clearpage
\noindent\mbox{}\hfill THE EARTH ELLIPSOID.\hfill 69\par
\[
\begin{array}{r@{\quad}c@{\quad}r}
10\,016\,491^{\mathrm m}{,}8250&&\\
47{,}5954&&\\
0{,}1682&&\\\cline{1-1}
10\,016\,539^{\mathrm m}{,}5886&\ldotp\,\ldotp\,\ldotp\,\ldotp\,\ldotp\,\ldotp\,\ldotp\,\ldotp\,\ldotp\,\ldotp\,\ldotp\,\ldotp\,\ldotp\,\ldotp&7{,}000\,7177\,118\\
&\left[\dfrac\pi2\right]\ldots&0{,}196\,1198\,770\\\cline{3-3}
\multicolumn{2}{r}{\log a\text{ (in metres)}=}&6{,}804\,5978\,348
\end{array}
\]
[Given the] meridian quadrant $=5130878{,}3$ toises[, one has]
\[
\log a\text{ (in toises)}=6{,}514\,7893\,106.
\]
\begin{center}\rule{.25\textwidth}{.4pt}\end{center}
Numbers concerning the Earth's spheroid.
\[
\text{Flattening}=\frac1{302{,}78},\quad ee=\frac{604{,}56}{(302{,}78)^2},\quad\log ee=7{,}81918\,50399{\scriptstyle45}.
\]
[The circumference of the meridian $=$] the circumference of the equator times
\[
\begin{aligned}
\bigg(1-\tfrac14ee&-\frac{1.3}{4.16}e^4-\frac{1.3.15}{4.16.36}e^6-\frac{1.3.15.35}{4.16.36.64}e^8-\frac{1.3.15.35.63}{4.16.36.64.100}e^{10}\\
&\bigg[-\frac{1.3.15.35.63.99}{4.16.36.64.100.144}e^{12}-\frac{1.3.15.35.63.99.143}{4.16.36.64.100.144.196}e^{14}\bigg]-\ldots\bigg).
\end{aligned}
\]
\[
\begin{array}{rcl}
\tfrac14ee&=&\multicolumn{1}{r}{0{,}00164\,86370\,23537\,4099}\\[3pt]
\tfrac3{64}e^4&=&\multicolumn{1}{r}{20385\,03026\,5312}\\[3pt]
\tfrac5{256}e^6&=&\multicolumn{1}{r}{56\,01252\,6094}\\[3pt]
\tfrac{175}{16384}e^8&=&\multicolumn{1}{r}{20200\,3209}\\[3pt]
\tfrac{441}{65536}e^{10}&=&\multicolumn{1}{r}{83\,9236}\\[3pt]
\tfrac{4851}{1048576}e^{12}&=&\multicolumn{1}{r}{3805}\\[3pt]
\left[\tfrac{693693}{205520896}e^{14}\right]&=&\multicolumn{1}{r}{18}\\\cline{3-3}
&&0{,}00165\,06811\,48101\,1773\\[3pt]
\text{Sum of the series}&=&0{,}99834\,93188\,51898\,8227
\end{array}
\]
\[
\begin{array}{rcl}
\text{Logarithm.}&\ldots&9{,}99928\,25259{\scriptstyle92}\\
\dfrac{20\,000\,000}{\pi}&\ldotp\,\ldotp\,\ldotp\,\ldotp\,\ldotp\,\ldotp&6{,}80388\,01229{\scriptstyle70}\\\cline{3-3}
\text{\raisebox{.5ex}{.}}\quad\log a&=&6{,}80459\,75969{\scriptstyle78}.
\end{array}
\]

\clearpage
\noindent 70\hfill POSTHUMOUS PAPERS.\hfill\mbox{}\par
\begin{center}[2.]\par
[Equation of the vertical plane of the ellipsoid of revolution.]\end{center}
Consider a plane that is perpendicular to the surface of the ellipsoid at longitude $0$ and latitude
$\varphi^0$ and that intersects the northern part of the meridian there at an angle
$\zeta^0$. For points in the intersection of this plane and the surface, the condition is
\[
\frac{\cos\varphi\cos\lambda\sin\varphi^0}{\sqrt{(1-ee\sin\varphi^2)}}+\frac{\cos\varphi\sin\lambda\operatorname{cotang}\zeta^0}{\sqrt{(1-ee\sin\varphi^2)}}-\frac{(1-ee)\sin\varphi\cos\varphi^0}{\sqrt{(1-ee\sin\varphi^2)}}=\frac{ee\cos\varphi^0\sin\varphi^0}{\sqrt{(1-ee\sin\varphi^{0^2})}}
\]
or, more simply,
\[
\cos\varphi\cos\lambda\sin\varphi^0+\cos\varphi\sin\lambda\operatorname{cotang}\zeta^0-(1-ee)\sin\varphi\cos\varphi^0=ee\cos\varphi^0\sin\varphi^0\sqrt{\frac{1-ee\sin\varphi^2}{1-ee\sin\varphi^{0^2}}}.
\]
\begin{center}[3.]\end{center}
Equation of the ellipsoid [of revolution] with respect to a tangent plane:
\[
xx(1-ee\sin\varphi^2)+yy(1-ee)+xzee\sin2\varphi+zz(1-ee\cos\varphi^2)-\frac{2bbz}{a\sqrt{(1-ee\sin\varphi^2)}}=0.
\]
[The $x$-axis is tangent to the meridian ellipse, positive towards the south. The $z$-axis
coincides with the direction of the normal and is positive towards the interior of the ellipsoid.]

[One also has:]
\[
\left(\frac{x\sin\varphi-z\cos\varphi}{a}\right)^2+\left(\frac{x\cos\varphi+z\sin\varphi}{b}\right)^2+\frac{yy}{aa}=\frac{2z}{a\sqrt{(1-ee\sin\varphi^2)}}.
\]
[If one sets]
\[
\begin{aligned}
x&=s\cos\delta\cos\theta\\
y&=s\cos\delta\sin\theta\\
z&=s\sin\delta,
\end{aligned}
\]
[where $\delta$ is the angle of depression and $\theta$ the azimuth, the result is]
\[
s\left(1+\frac{ee}{1-ee}(\cos\varphi\cos\theta\cos\delta+\sin\varphi\sin\delta)\right)^2\sqrt{(1-ee\sin\varphi^2)}=2a\sin\delta.
\]
[Approximately one has]
\[
\sin\lambda=\frac{s\cos\delta\sin\theta\sqrt{(1-ee\sin\varphi^2)}}{a\cos\varphi}.
\]
\begin{center}\rule{.18\textwidth}{.4pt}\end{center}

\clearpage
\noindent\mbox{}\hfill THE EARTH ELLIPSOID.\hfill 71\par
\vspace{12mm}
\begin{center}\textbf{REMARKS.}\end{center}
{\small
Note [1] is taken from a notebook. The corresponding numerical values were found in two other
notebooks. The numbers concerning the ellipsoid are entered on the last page of \textsc{Gauss}'s
copy of »Mathematischen Abhandlungen von Dr.\ H.\ F.\ \textsc{Scherk}. Berlin, 1825«. The flattening
$\tfrac1{310}$ was the value used by General von \textsc{Müffling} as the basis for the Prussian
surveys. As the letter to \textsc{Olbers} of 18 April 1822 shows, at the beginning of his arc
measurement work \textsc{Gauss} inadvertently used the flattening $\dfrac1{302{,}68}$ in place of
\textsc{Walbeck}'s value. Indeed, several tables and calculations are based on it. \textsc{Gauss}
at first intended to keep this value, but later decided (see \textsc{Gauss}'s letter to
\textsc{Olbers} of 1 March 1827) to adopt \textsc{Walbeck}'s value $\dfrac1{302{,}78}$. This is
therefore the basis for the calculation of \textsc{Gauss}'s measurements. The length of the meridian
quadrant according to \textsc{Walbeck} is $5130878{,}22$ toises (Dissertatio de forma et magnitudine
telluris ex dimensis arcubus meridiani definiendis. Aboae, 1819).

[2] and the first part of [3] belong to the same notebook as [1]. The second part of [3]
was found on a single leaf.
\par\hfill\textsc{Krüger, Börsch.}\par
}
\begin{center}\rule{.25\textwidth}{.4pt}\end{center}
